\documentclass[a4,11pt]{aleph-notas}

% -- Paquetes adicionales
\usepackage{enumitem}
\usepackage{aleph-comandos}

% -- Datos
\institucion{Facultad de Ciencias Exactas, Naturales y Ambientales}
\carrera{Catálogo STEM}
\asignatura{Matemática III}
\tema{Ejercicios 02: Geometría de curvas}
\autor{Andrés Merino}
\fecha{Periodo 2026-1}

\logouno[0.16\textwidth]{Logos/logoPUCE_04_ac}
\definecolor{colortext}{HTML}{1957d1}
\definecolor{colordef}{HTML}{1957d1}
\fuente{montserrat}

% -- Comandos adicionales
\setlist[enumerate]{label=\roman*.}

\begin{document}

\encabezado

\begin{ejer}
En cada uno de los literales siguientes, $\alpha(t)$ representa el vector posición en el instante $t$ correspondiente a una partícula que se mueve sobre una curva $C$. En cada caso, determinar los vectores velocidad $v(t)$ y aceleración $a(t)$, para cada $t\in\R$.	
    \begin{enumerate}
    \item 
        \[
            \funcion{\alpha}{\R}{\R^3}{t}{(\cos(t),\sen(t),e^t)}
        \]
    \item
        \[
            \funcion{\alpha}{\R}{\R^3}{t}{t\veci+\sen(t)\vecj+(1-\cos(t))\veck}
        \]
    \item
    	\[
        	\funcion{\alpha}{\R}{\R^3}{t}{(t-\sen(t))\veci+(1-\cos(t))\vecj+4\sen\l(\frac{t}{2}\r)\veck}
       	\]
    \end{enumerate}
\end{ejer}

\begin{proof}[Solución]
\begin{enumerate}
    \item
        Tenemos que
        \[
            v(t)=\alpha'(t)=(-\sen(t),\cos(t),e^t)
            \texty
            a(t)=\alpha''(t)=(-\cos(t),-\sen(t),e^t).
        \]
    \item
        Tenemos que
        \[
            v(t)=\alpha'(t)=\veci +\cos(t)\vecj +\sen(t)\veck
            \texty
            a(t)=\alpha''(t)= -\sin(t)\vecj +\cos(t)\veck.
        \]
    \item
    	Tenemos que
        \[
            v(t)=\alpha'(t)=(1-\cos(t))\veci +\sen(t)\vecj +2\cos\l(\frac{t}{2}\r)\veck
        \]
        y
        \[
            a(t)=\alpha''(t)= \sen(t)\veci +\cos(t)\vecj -\sin\l(\frac{t}{2}\r)\veck. \qedhere
        \]
    \end{enumerate}
\end{proof}

%%%%%%%%%%%%%%%%%%%%%%%%%%%%%%%%%%%%%%%%%%%
\begin{ejer}
Dada la trayectoria
    \[
    	\funcion{\alpha}{\R}{\R^3}{t}{3t\cos(t) \ \veci+3t\sen(t)\   \vecj\ + 4t\ \veck,}
    \]
    la cual representa el vector posición de un movimiento, determinar los vectores velocidad $v(t)$ y aceleración $a(t)$ y la rapidez $\rho(t)$, para cada $t\in\R$.
\end{ejer}

%%%%%%%%%%%%%%%%%%%%%%%%%%%%%%%%%%%%%%%%%%%
\begin{ejer}
Considere las trayectorias $\func{\alpha,\beta}{\l[0,+\infty\r[}{\R^3}$, definidas por $\alpha(t)=(\cos(t),\sen(t),t)$ y $\beta(t)=(1,0,t)$ para todo $t\in\l[0,+\infty\r[$. Estas trayectorias representan los vectores posición en un instante $t$ de dos partículas que se mueven sobre dos curvas descritas por $\alpha$ y $\beta$ respectivamente. ¿En qué puntos, si los hay, estas partículas se encuentran?
\end{ejer}

\begin{proof}[Solución]
Buscamos los elementos $t\in \l[0,+\infty\r[$ tal que
    \[
        \alpha(t)=\beta(t),
    \]
    es decir,
    \[
       (\cos(t),\sen(t),t)=(1,0,t).
    \]
    Así, se tienen las ecuaciones
    \[
        \cos(t)=1,
        \qquad
        \sen(t)=0
        \texty
        t=t.
    \]
    Por lo tanto, se tiene que las partículas se encuentran cuando
    \[
        t\in\{2k\pi:k\in\N\},
    \]
    y lo hacen en los puntos 
    \[
        \{(1,0,2k\pi):k\in\N\}.\qedhere
    \]
\end{proof}

%%%%%%%%%%%%%%%%%%%%%%%%%%%%%%%%%%%%%%%%%%%
\begin{ejer}
Suponga que no conocemos la trayectoria de un planeador, solo su vector de
aceleración, que está dada por $a(t)= \l(-3 \cos t, -3 \sen t, 2 \r)$, para $t\in \l[0,+\infty\r[$. También sabemos que inicialmente
(en el instante $t=0$), el planeador partió del punto $(3, 0, 0)$ con una velocidad inicial $v(0)=\l(0,3,0\r)$. Obtenga la posición del planeador como una función de $t$.
\end{ejer}

\begin{proof}[Solución]
Sea $\func{r}{\l[0,+\infty\r[}{\R^3}$ la posición de la partícula. Conocemos que la velocidad $v$ es la derivada de la posición $r$ y la aceleración $a$, la derivada de $v$. Entonces, teniendo la función aceleración, la integramos dos veces para obtener la posición, entonces
    \[
        a(t)=\dfrac{dv}{dt}(t)=\dfrac{d}{dt}\dfrac{dr}{dt}(t)= \l(-3 \cos t, -3 \sen t, 2 \r)
    \]
    para $t\in\l[0,+\infty\r[$, con condiciones iniciales:
    \[
        v(0)= \l(0,3,0\r)
        \texty
        r(0)= \l(3,0,0\r).
    \]
    Al integrar ambos lados de esta expresión tenemos
    \[
        v(t)= \l(-3 \sen t+C_1, 3 \cos t+C_2, 2t +C_3 \r),
    \]
    utilizando la condición inicial $v(0)= \l(0,3,0\r)$ se determina que $C_1=0$, $C_3=0$ y $C_3=0$.

    Luego, la velocidad del planeador es
    \[
        v(t)= \l(-3 \sen t, 3 \cos t, 2t \r)
    \]
    para todo $t\in\l[0,+\infty\r[$. Al integrar la velocidad $v$ se obtiene la posición del planeador
    \[
        r(t)= \l(3 \cos t + C_4, 3 \sen t + C_5, t^2 + C_6 \r),
    \]
    utilizando la condición inicial $r(0)= \l(3,0,0\r)$ de donde, $C_4=0$, $C_5=0$ y $C_6=0$.

    Por lo tanto la posición del planeador es
    \[
        r(t)= \l(3 \cos t, 3 \sen t, t^2\r)
    \]
    para todo $t\in\l[0,+\infty\r[$.
\end{proof}

%%%%%%%%%%%%%%%%%%%%%%%%%%%%%%%%%%%%%%%%%%%
\begin{ejer}
Un cohete descompuesto se mueve de acuerdo con la trayectoria
\[
\funcion{\alpha}{[0,+\infty[\setminus \lbrace 0 \rbrace}{\R^3}
{t}{\alpha(t)=\l(e^{2t},3t^3-2t,t-\dfrac{1}{t}\r)}
\]
con la esperanza de llegar a una estación de reparación ubicada en el punto $(7e^4,35,5)$. El tiempo se mide en minutos y las coordenadas espaciales se miden en millas. Los motores del cohete se apagan de súbito a los dos minutos, si sobre el cohete no actúa ninguna fuerza externa, ¿llegará el cohete con su solo impulso a la estación de reparación?
\end{ejer}

\begin{proof}[Solución]
Notemos que $\alpha$ es derivable con derivada
\[
\alpha'(t) = \l(2e^{2t},9t^2-2,1+\dfrac{1}{t^2}\r)
\]
para cada $t\in \l] 0,\infty \r[ $. Entonces, a los dos minutos, el cohete se encuentra en la posición $\alpha(2)=\l(e^4,20,\dfrac{3}{2}\r)$ con velocidad $\alpha'(2)=\l(2e^4,34,\dfrac{5}{4}\r)$. 

Ya que el cohete deja de funcionar a los dos minutos, este deberá continuar por la trayectoria de la recta tangente a $\alpha$ en 2, es decir, la recta
\begin{align*}
    L(\alpha(2),\alpha'(2)) &= \l\{\l(e^4,20,\dfrac{3}{2}\r) + (t-2)\l(2e^4,34,\dfrac{5}{4}\r)\,:\,t \geq 2\r\}\\
    &= \l\{\l((2t-3)e^4, 34t-48,\dfrac{5}{4}t-1\r)\,:\,t \geq 2\r\}\\
\end{align*}
Si el cohete logra alcanzar la estación de reparación debe tenerse que $(7e^4,35,5)\in L(\alpha(2),\alpha'(2))$, es decir, se tiene
\[
(7e^4,35,5) = \l((2t-3)e^4, 34t-48,\dfrac{5}{4}t-1\r)
\]
para algún $t \geq 2$. En particular tenemos
\[
7e^4 = (2t-3)e^4
\]
lo que se verifica cuando $t=5$. Notemos ahora que si $t=5$ el cohete se encontrará en la posición $\l(7e^4,122,\dfrac{21}{4}\r)\neq(7e^4,35,5)$.

Por lo tanto, el cohete no logrará llegar hasta la estación de reparación.
\end{proof}

%%%%%%%%%%%%%%%%%%%%%%%%%%%%%%%%%%%%%%%%%%%
\begin{ejer}
Sea la trayectoria 
    \[
        \funcion{\alpha}{[-1,3]}{\R^2}{t}{(t,\sen(\pi t)).}
    \]
    Si $\alpha$ describe la trayectoria de una mosca, halle la recta tangente a su trayectoria en $t=1$. Interprete este resultado. Encuentre el desplazamiento total de la mosca.
\end{ejer}

%%%%%%%%%%%%%%%%%%%%%%%%%%%%%%%%%%%%%%%%%%%
\begin{ejer}
Dado un  vector fijo no nulo $b\in\R^3$, se tiene un movimiento con vector posición definido por $\func{\alpha}{\R}{\R^3}$, que satisface
    \[
        \alpha'(t) = b\times \alpha(t)
    \]
    para todo $t\in\R$. Demostrar que el vector aceleración es siempre ortogonal a $b$.
\end{ejer}

\begin{proof}[Solución]
Calculemos el vector aceleración para $t\in\R$,
    \[
        a(t)=\alpha''(t)=(b\times \alpha(t))'=0+b\times \alpha'(t)=b\times \alpha'(t).
    \]
    Ahora, para comprobar que es ortogonal a $b$, notemos que
    \[
        b\cdot a(t) = b\cdot (b\times \alpha'(t)) = 0
    \]
    para todo $t\in\R$, así, el vector aceleración es siempre ortogonal a $b$.
\end{proof}

%%%%%%%%%%%%%%%%%%%%%%%%%%%%%%%%%%%%%%%%%%%
\begin{ejer}
Dada la trayectoria
	\[
        \funcion{\alpha}{\R}{\R^3}{(t)}{b\cos(\omega t)\veci+b\sen(\omega t)\vecj +(c\omega t)\veck}
    \]
    donde $\omega$ es una constante positiva y $b,c\in \R$. Compruebe que los vectores velocidad $v(t)$ y aceleración $a(t)$ tienen longitud constante para todo $t\in \R$ y que
    \[
        \frac{\|v(t)\times a(t)\|}{\|v(t)\|^3}=\frac{b}{b^2+c^2},
    \]
    para todo $t\in \R$.
\end{ejer}

%%%%%%%%%%%%%%%%%%%%%%%%%%%%%%%%%%%%%%%%%%%
\begin{ejer}
Compruebe que, para cualquier movimiento, el producto escalar de los vectores velocidad y aceleración es igual a la mitad de la derivada del cuadrado de la rapidez, es decir, dado un movimiento definido por la trayectoria $\func{\alpha}{I}{\R^n}$, se tiene que
	\[
        v(t)\cdot a(t)=\dfrac{1}{2} \dfrac{d}{dt} (\rho^2(t))
    \]
    para todo $t\in\R$.
\end{ejer}

\begin{proof}[Solución]
Se conoce que $\|v(t)\|=\rho(t)$, para todo $t\in I$, por lo tanto
    \begin{align*}
        \rho^2(t) = v(t)\cdot v(t).
	\end{align*}
    Derivando esto respecto al tiempo, tenemos
    \begin{align*}
		\dfrac{d}{dt} (\rho^2(t))
		    & =\dfrac{d}{dt} (v(t) \cdot v(t)) \\
            & = v'(t) \cdot v(t) + v(t) \cdot v(t)'\\
            & = 2 v(t) \cdot v(t)',
	\end{align*}
	para todo $t\in I$. Por lo tanto, como la aceleración es la derivada del vector velocidad, 
	\[
        \dfrac{d}{dt} (\rho^2(t))=2v(t)\cdot a(t),
   	\]
	para todo $t\in I$; reordenamos la expresión y obtenemos
	\[
   		\dfrac{1}{2} \dfrac{d}{dt} (\rho^2(t))=v(t)\cdot a(t),
    \]
    para todo $t\in I$, como queríamos.
\end{proof}

%%%%%%%%%%%%%%%%%%%%%%%%%%%%%%%%%%%%%%%%%%%
\begin{ejer}
Dada la trayectoria 
	\[	
        \funcion{\alpha}{\l[0,+\infty\r[}{\R^n}{t}{\alpha(t)=c+\dfrac{1-e^{-kt}}{k}d,}
    \]
	donde $k>0$ y $c,d\in\R^n$.	Muestre que $\alpha$ verifica
    \[
        \alpha''=-k\alpha'.
    \]
	Hallar la relación de $c$ y $d$ con la posición inicial y calcular
    \[
    	\lim_{t\to+\infty} (\alpha(t) - \alpha(0))
    \]
    en términos de la velocidad inicial (a este límite se lo conoce como la velocidad total de desplazamiento de la partícula).
\end{ejer}

\begin{proof}[Solución]
Es fácil notar que $\alpha$ es infinitamente derivable y se tiene que
	\[
		\alpha'(t)=e^{-kt}d 
		\texty
		\alpha''(t)=-ke^{-kt}d
	\]
	para cada $t\in\R$. Así,
	\[
		\alpha''=-k\alpha'.
	\]
	Por otro lado, puesto que
	\[
		\alpha(0)=c\quad \text y \quad 
		\alpha'(0)=d,
	\]
	es decir, $c$ es la posición inicial y $d$ es la velocidad inicial. Además, se tiene que
	\[
		\lim_{t\to+\infty}( \alpha(t) - \alpha(0) )
		=\lim_{t\to+\infty} \l(c+\dfrac{1-e^{-kt}}{k}d\r) - c
		=\lim_{t\to+\infty} \dfrac{1-e^{-kt}}{k}d
		=\dfrac{1}{k}d,
	\]
	dado que $k>0$. Finalmente, la velocidad total de desplazamiento de la partícula es
	\[
		\lim_{t\to+\infty}( \alpha(t) - \alpha(0) )
		=\dfrac{1}{k}\alpha'(0). \qedhere
	\]
\end{proof}

%%%%%%%%%%%%%%%%%%%%%%%%%%%%%%%%%%%%%%%%%%%
\begin{ejer}
Compruebe que, en un movimiento representado por la función $\func{\alpha}{I}{\R^3}$, si el vector aceleración $a(t)$ es cero para todo $t\in I$, el movimiento es rectilíneo.
\end{ejer}

\begin{proof}[Solución]
Tenemos que
    \[
    	\funcion{\alpha}{I}{\R^3}{t}{(\alpha_1(t),\alpha_2(t),\alpha_3(t)),}
    \]
    y dado que $a(t)=0$ para todo $t\in I$, se tiene que
    \[
        \alpha_1''(t)=0,
        \qquad
        \alpha_2''(t)=0
        \texty
        \alpha_3''(t)=0,
    \]
    y por lo tanto
    \[
        \alpha_1(t)=a_1+tb_1,
        \qquad
        \alpha_2(t)=a_2+tb_2
        \texty
        \alpha_3(t)=a_3+tb_3,
    \]
    donde $a_i,b_i\in\R$ para todo $i\in\{1,2,3\}$.
    Así, el vector posición es
    \[
        \alpha(t)=(a_1+tb_1,a_2+tb_2,a_3+tb_3)=a+tb
    \]
    para todo $t\in I$, donde $a=(a_1,a_2,a_3)$ y $b=(b_1,b_2,b_3)$, y por lo tanto el movimiento descrito por $r$ es rectilíneo.
\end{proof}

%%%%%%%%%%%%%%%%%%%%%%%%%%%%%%%%%%%%%%%%%%%
\begin{ejer}
Verificar que la trayectoria 
    \[
        \funcion{\alpha}{\R}{\R^3}{t}{(\sen t^2,\cos t^2,e)}
    \]
    describe un movimiento planar.
\end{ejer}

\begin{proof}[Solución]
Para esto se debe verificar que la dirección del vector binormal es constante, entonces
    \begin{align*}
        \alpha'(t)&=(2t\cos t^2,-2t\sen t^2,0), \\
        \rho(t)&=2t
    \end{align*}
    de donde 
    \[
        T(t)=(\cos t^2,-\sen t^2,0)
    \]
    y
    \begin{align*}
        T'(t)&=(-2t\sen t^2,-2t\cos t^2,0)\\
        \|T'(t)\|&=2t\\
        N(t)&=(-\sen t^2,-\cos t^2,0).
    \end{align*}
    Finalmente
    \begin{align*}
        B(t)&=T(t) \times N(t)\\
        &=(\cos t^2,-\sen t^2,0) \times (-\sen t^2,-\cos t^2,0)\\
        &=(0,0,-1)
    \end{align*}
    con lo que se verifica que el movimiento descrito es planar.
\end{proof}

%%%%%%%%%%%%%%%%%%%%%%%%%%%%%%%%%%%%%%%%%%%
\begin{ejer}
Dada una curva $C$ en $\R^3$, con dos parametrizaciones equivalente
    \[
        \func{\alpha}{I}{\R^3}
        \texty
        \func{\beta}{J}{\R^3}
    \]
    con cambio de parámetro $\func{u}{I}{J}$ creciente, demostrar que los vectores principales no se ven afectados por el cambio de parámetro.
\end{ejer}

\begin{proof}[Solución]
Dado que $u$ es el cambio de parámetro, se tiene que
    \[
        \beta(u(t))=\alpha(t)
    \]
    para todo $t\in I$. Además, dado que $u$ es creciente, se tiene que $u'(t)\geq 0$ para todo $t\in\R$.
    
    Fijemos $t\in \R$ y llamemos $\tau=u(t)$. Calculemos los vectores principales en en el punto $\beta(\tau)$. Con la parametrización $\beta$, tenemos que
    \[
        T_\beta(\tau)= \frac{\beta'(\tau)}{\|\beta'(\tau)\|},
        \qquad
        N_\beta(\tau)= \frac{T_\beta'(\tau)}{\|T_\beta'(\tau)\|}
        \texty
        B_\beta(\tau)= T_\beta(\tau)\times N_\beta(\tau).
    \]
    Ahora, hallemos los vectores en el mismo punto pero con la parametrización $\alpha$, es decir en el punto $\beta(\tau)=\beta(u(t))=\alpha(t)$ y tenemos
    \[
        T_\alpha(t)= \frac{\alpha'(t)}{\|\alpha'(t)\|},
        \qquad
        N_\alpha(t)= \frac{T_\alpha'(\tau)}{\|T_\alpha'(t)\|}
        \texty
        B_\alpha(t)= T_\alpha(t)\times N_\alpha(t).
    \]
    Pero notemos que
    \begin{align*}
        T_\alpha(t)
            & = \frac{\alpha'(t)}{\|\alpha'(t)\|} = \frac{(\beta(u(t)))'}{\|(\beta(u(t)))'\|}\\[1mm]
            & = \frac{u'(t)\beta'(u(t))}{\|u'(t)\beta'(u(t))\|} = \frac{u'(t)\beta'(u(t))}{|u'(t)|\|\beta'(u(t))\|}\\[1mm]
            & = \frac{\beta'(u(t))}{\|\beta'(u(t))\|} = \frac{\beta'(\tau)}{\|\beta'(\tau)\|}\\[1mm]
            & = T_\beta(\tau),
    \end{align*}
    es decir, $T\alpha(t)=T_\beta(u(t))$. Por otro lado, 
    \begin{align*}
        N_\alpha(t)
            & = \frac{T_\alpha'(t)}{\|T_\alpha'(t)\|} = \frac{(T_\beta(u(t)))'}{\|(T_\beta(u(t)))'\|}\\[1mm]
            & = \frac{u'(t)T_\beta'(u(t))}{\|u'(t)T_\beta'(u(t))\|} = \frac{u'(t)T_\beta'(u(t))}{|\beta'(t)|\|T_\beta'(u(t))\|}\\[1mm]
            & = \frac{T_\beta'(u(t))}{\|T_\beta'(u(t))\|} = \frac{T_\beta'(\tau)}{\|T_\beta'(\tau)\|}\\[1mm]
            &= N_\beta(\tau).
    \end{align*}
    Finalmente
    \[
        B_\alpha(t)= T_\alpha(t)\times N_\alpha(t) 
        = T_\beta(\tau)\times N_\beta(\tau)=B_\beta(\tau),
    \]
    es decir,
    \[
        T_\alpha(t)= T_\beta(\tau),
        \qquad
        N_\alpha(t)= N_\beta(\tau)
        \texty
        B_\alpha(t)= B_\beta(\tau).\qedhere
    \]
\end{proof}

%%%%%%%%%%%%%%%%%%%%%%%%%%%%%%%%%%%%%%%%%%%
\begin{ejer}
Demostrar que el módulo de la componente normal del vector aceleración en $t$ es $\dfrac{\|v(t) \times a(t)\|}{\|v(t)\|}$ para cualquier curva.
\end{ejer}

\begin{proof}[Solución]
Sabemos que $a(t)$ se puede escribir como una combinación lineal de los vectores tangente y normal, es decir,
    \[
        a(t)=\alpha T(t)+\beta N(t),
    \]
    donde $\alpha,\beta\in\R$.
    Ahora, tenemos que
    \[
        v(t)\times a(t)=\alpha v(t)\times T(t)+\beta v(t)\times N(t).
    \]
    Notemos que, gracias a que $v(t)$ y $T(t)$ son paralelos, $v(t)\times T(t)=0$, por lo tanto
    \[
        v(t)\times a(t)=\beta v(t)\times N(t).
    \]
    Así, se tiene que
    \[
         \| v(t)\times a(t) \| = |\beta| \| v(t)\times N(t) \|,
    \]
    y por lo tanto
    \begin{align*}
        |\beta| 
            & = \frac{\| v(t)\times a(t) \|}{\| v(t)\times N(t) \|} \\
            & = \frac{\| v(t)\times a(t) \|}{\| v(t) \| \| N(t) \| \sen\l(\frac{\pi}{2}\r)} \\
            & = \frac{\| v(t)\times a(t) \|}{\| v(t) \| \| N(t) \|}\\
            & = \frac{\| v(t)\times a(t) \|}{\| v(t) \|}.
    \end{align*}
    Luego, la componente normal de la aceleración es
    \[
        \frac{\| v(t)\times a(t) \|}{\| v(t) \|} N(t)
    \]
\end{proof}

%%%%%%%%%%%%%%%%%%%%%%%%%%%%%%%%%%%%%%%%%%%
\begin{ejer}
Dado el movimiento representado por 
	\[
		\funcion{\alpha}{\R}{\R^3}{t}{(\sen(t),\cos(t),t^2)}.
	\]
	Verifique que se cumple que
	\[
        a(t)=\rho'(t) T(t)+\rho(t) T' (t).
	\]
	para cada $t\in\R$.
\end{ejer}

\begin{proof}[Solución]
Se tiene que
	\begin{align*}
        v(t)
        &=\alpha'(t)=(\cos(t),-\sen(t),2t),
        &\rho(t)&=\|v(t)\|=\sqrt{1+4t^2},\\
        a(t)&=\alpha''(t)=(-\sen(t),-\cos(t),2),
        &\rho'(t)&=\dfrac{4t}{\sqrt{1+4t^2}}
	\end{align*}
	para cada $t\in\R$. Puesto que $\alpha'(t)\neq0$ para cada $t\in\R$, se sigue que
	\begin{align*}
		T(t)=\l(\dfrac{\cos(t)}{\sqrt{1+4t^2}},\dfrac{-\sen(t)}{\sqrt{1+4t^2}},\dfrac{2t}{\sqrt{1+4t^2}}\r)
	\end{align*}
	y además, 
	\[
		T'(t)
		=\l(\dfrac{-(1+4t^2) \sen(t) -4t\cos(t))}{(1 + 4 t^2)^\frac{3}{2}},
		\dfrac{ 4t\sen(t)-(1+4t^2) \cos(t))}{(1 + 4 t^2)^\frac{3}{2}},
		\dfrac{ 2}{(1 + 4 t^2)^\frac{3}{2}}\r)
	\]
	para todo $t\in\R$. Finalmente, para cada $t\in\R$, se tiene que
	\begin{align*}
		\rho'(t) T(t)+&\rho(t) T' (t)=\\
		&= \dfrac{4t}{\sqrt{1+4t^2}}\l(\dfrac{\cos(t)}{\sqrt{1+4t^2}},\dfrac{-\sen(t)}{\sqrt{1+4t^2}},\dfrac{2t}{\sqrt{1+4t^2}}\r)\\
	    &\phantom{=}+\sqrt{1+4t^2}\l(\dfrac{-(1+4t^2) \sen(t) -4t\cos(t))}{(1 + 4t^2)^\frac{3}{2}},
			\dfrac{ 4t\sen(t)-(1+4t^2) \cos(t))}{(1 + 4 t^2)^\frac{3}{2}},
			\dfrac{ 2}{(1 + 4 t^2)^\frac{3}{2}}\r)\\[2mm]
		&= 4t\l(\dfrac{\cos(t)}{1+4t^2},\dfrac{-\sen(t)}{1+4t^2},\dfrac{2t}{1+4t^2}\r)\\
		&\phantom{=}+\l(\dfrac{-(1+4t^2) \sen(t) -4t\cos(t))}{1 + 4t^2},
			\dfrac{ 4t\sen(t)-(1+4t^2) \cos(t))}{1 + 4t^2},
			\dfrac{ 2}{1 + 4t^2}\r)\\[2mm]
		&=\l(\dfrac{-(1+4t^2) \sen(t))}{1 + 4t^2},
			\dfrac{-(1+4t^2) \cos(t))}{1 + 4t^2},
			\dfrac{8t^2+2}{1 + 4t^2}\r)\\[1mm]
		&=(-\sen(t),-\cos(t),2)\\
		&=a(t).
	\end{align*}
	Así, se concluye que
	\[
		a(t)=\rho'(t) T(t)+\rho(t) T' (t)
	\]
	para cada $t\in\R$.
\end{proof}

%%%%%%%%%%%%%%%%%%%%%%%%%%%%%%%%%%%%%%%%%%%
\begin{ejer}
Calcule la longitud de $C$ representada por la trayectoria $\func{\alpha}{[0,2\pi]}{\R^3}$ donde $\alpha$ se define como en el ejercicio anterior. Además, calcular su vector curvatura y su curvatura.
\end{ejer}

\begin{proof}[Solución]
Utilizando los cálculos anteriores y la definición, se tiene que
	\[
		\ell(C)=\int_0^{2\pi}\rho(t)dt=\int_0^{2\pi}\sqrt{1 + 4t^2}dt=
		\pi \sqrt{1 + 16 \pi^2} + \dfrac{1}{4} \senh^{-1}(4\pi)\approx 40.4.
	\]
	Por otro lado, se tiene que
	\begin{align*}
		K(t)=\dfrac{T'(t)}{\rho(t)}
		&=\dfrac{1}{\sqrt{1+4t^2}}\l(\dfrac{-(1+4t^2) \sen(t) -4t\cos(t))}{(1 + 4 t^2)^\frac{3}{2}},
		\dfrac{ 4t\sen(t)-(1+4t^2) \cos(t))}{(1 + 4 t^2)^\frac{3}{2}},
		\dfrac{ 2}{(1 + 4 t^2)^\frac{3}{2}}\r)\\[2mm]
		&=\l(\dfrac{-(1+4t^2) \sen(t) -4t\cos(t))}{(1 + 4 t^2)^2},
		\dfrac{ 4t\sen(t)-(1+4t^2) \cos(t))}{(1 + 4 t^2)^2},
		\dfrac{ 2}{(1 + 4 t^2)^2}\r)
	\end{align*}
	y
	\begin{align*}
		\kappa(t)=\dfrac{\|T'(t)\|}{\rho(t)}
		&=\dfrac{1}{(1+4t^2)^2}
		\sqrt{(-(1+4t^2) \sen(t) -4t\cos(t)))^2
				+(4t\sen(t)-(1+4t^2) \cos(t)))^2
				+4}\\
		&=\dfrac{1}{(1+4t^2)^2}
			\sqrt{
			(1+4t^2)^2 +16t^2+4}\\
		&=\dfrac{1}{(1+4t^2)^2}
			\sqrt{
				(4 t^4 + 16 t^2 + 5}
	\end{align*}
	para todo $t\in\lcl0,2\pi\rcl$.
\end{proof}

%%%%%%%%%%%%%%%%%%%%%%%%%%%%%%%%%%%%%%%%%%%
\begin{ejer}
Dada la trayectoria 
 	\[
        \funcion{\alpha}{\left[0,2\pi \right]}{\R^2}{t}{a(1-\cos (t)) \veci+a(t-\sen (t)) \vecj}
    \]
	con $a>0$. Hallar 
	\begin{enumerate}
	    \item La longitud de la curva $C$ descrita por $\alpha$. (Distancia) 
	    \item Desplazamiento entre los puntos inicial y final de la curva $C$.
	\end{enumerate}
\end{ejer}

\begin{proof}[Solución]
\begin{enumerate}
	  \item
	 	El vector velocidad correspondiente es
    	\[
       		v(t)=\alpha'(t)= a \sen(t) \veci+(a-a \cos (t)) \vecj,
        \]
    	para todo $t\in[0,2\pi]$ y la rapidez asociada es
    	\[
            \rho(t) = \| v(t) \| = a \sqrt{2-2\cos(t)} 
      	\]
    	para todo $t\in\left[0,2\pi\right]$. Por la definición de la longitud de arco de la curva descrita por $r$, se tiene que
    	\begin{align*}
    		\ell(C)
                & =\int_0^{2\pi} \|v(t)\|\ dt\\
                & =\int_0^{2\pi}a \sqrt{2-2\cos(t)}\  dt\\
                & =8a.\qedhere
        \end{align*}
    \item
        El desplazamiento $\Delta r$ es la diferencia entre los puntos final e inicial, que son $\alpha(2\pi)$ y $\alpha(0)$ respectivamente, entonces
        \begin{align*}
            \Delta r&=a(0,2\pi)-a(0,0)\\
            &=a(0,2\pi)
        \end{align*}
        y su norma es $\| \Delta r\|=2a\pi$.
    \end{enumerate}
\end{proof}

%%%%%%%%%%%%%%%%%%%%%%%%%%%%%%%%%%%%%%%%%%%
\begin{ejer}
La trayectoria definida por
	\[
        \funcion{\alpha}{[0, 2\pi]}{\R^2}{t}
        {\alpha(t) = (a\sen{t}, b\cos{t})}
	\]
    donde $0 < b < a$ son números reales, es la parametrización de una elipse. Probar que la longitud $L$ de una elipse está dada por la integral
    \[
        L = 4a\int_{0}^{\pi/2}{\sqrt{1 - e^{2}\sen^{2}{t}}\ dt,}
    \]
    donde $e = \dfrac{\sqrt{a^{2} - b^{2}}}{a}$. Al número $e$ se lo llama excentricidad de la elipse.
\end{ejer}

%%%%%%%%%%%%%%%%%%%%%%%%%%%%%%%%%%%%%%%%%%%
\begin{ejer}
Dada la curva definida por la trayectoria
    \[
        \funcion{\alpha}{\R}{\R^3}
        {t}{\alpha(t) = (e^{t}\cos{t}, e^{t}\sen{t}, e^{t})}
	\]
    Hallar la curvatura y la ecuación del plano osculador cuando $t = 0$.
\end{ejer}

\begin{proof}[Solución]
Nótese que $\alpha$ es al menos dos veces derivable. En efecto, sus dos primeras derivadas $\alpha'$ y $\alpha''$ están dadas por
    \[
        \alpha'(t) = (e^{t}\cos{t} - e^{t}\sen{t}, e^{t}\cos{t} + e^{t}\sen{t}, e^{t}),
    \]
    y
    \[
        \alpha''(t) = (-2e^{t}\sen{t}, 2e^{t}\cos{t}, e^{t}),
    \]
    para todo $t \in \mathbb{R}$. Además,
    \[
        \| \alpha'(t) \| = \sqrt{(e^{t}\cos{t} - e^{t}\sen{t})^{2} + (e^{t}\cos{t} + e^{t}\sen{t})^2 + e^{2t}} = \sqrt{3} e^{t},
    \]
    para todo $t\in\R$. En particular, en $t = 0$, se tiene que $\alpha'(0)= (1,1,1)$, $\alpha''(0) = (0,2,1)$ y $\|\alpha'(0)\| = \sqrt{3}$, por tanto, la curvatura de $\alpha$ en $t = 0$, es
    \begin{align*}
        \kappa(0) 
            &= \frac{\|\alpha'(0) \times \alpha''(0) \|}{\| \alpha'(0) \|^{3}}\\
            &= \frac{\|(1,1,1) \times (0,2,1) \|}{\sqrt{3}^{3}}\\
            &= \frac{\|(-1,-1,2)\|}{3\sqrt{3}}\\
            &= \frac{\sqrt{2}}{3}.
    \end{align*}
    Para determinar la ecuación del plano osculador, calculemos el vector unitario tangente $T(t)$ y el vector normal principal $N(t)$. El vector unitario tangente está dado por
    \[
        T(t) = \frac{1}{\| \alpha'(t)\|}\alpha'(t) = \frac{\sqrt{3}}{3}\alpha'(t),
    \]
    para todo $t\in\R$ y cuya derivada $T'$ está dada por
    \[
        T'(t) = \frac{\sqrt{3}}{3}(-\cos{t} - \sen{t}, \cos{t} - \sen{t}, 0), 
    \]
    con norma $\| T'(t)\| = \frac{\sqrt{6}}{3}$, para todo $t \in\R$. El vector normal principal, está dado por
    \[
        N(t) = \frac{1}{\|T'(t)\|}T'(t) = \frac{\sqrt{2}}{2}T'(t).
    \]
    Nótese que un vector perpendicular al plano osculador es $B(t) = T(t) \times N(t)$ para todo $t\in\R$, por lo tanto, en $t = 0$, 
    \[
        B(0) = T(0) \times N(0) 
        = \l(\frac{\sqrt{3}}{3}\alpha'(0)\r)\times \l(\frac{\sqrt{2}}{2}T'(0)\r)
        =\frac{\sqrt{6}}{6} \l( (1,1,1) \times (-1, 1, 0) \r)
        = \frac{\sqrt{6}}{6} (-1, -1, 2).
    \]
    Con esto, se tiene que la ecuación del plano está dada por
    \[
        B(0)\cdot (x,y,z) = B(0)\cdot \alpha(0),
    \]
    es decir,
    \[
        (-1, -1, 2)\cdot(x,y,z) =(-1, -1, 2)\cdot(1, 0, 1)
    \]
    así, la ecuación del plano es
    \[
        -x - y + 2z = -1.\qedhere
    \]
\end{proof}

%%%%%%%%%%%%%%%%%%%%%%%%%%%%%%%%%%%%%%%%%%%
\begin{ejer}
Si un punto se mueve por una curva, parametrizada por la función
    \[
        \func{\alpha}{I}{\R^3,}
    \]
    de manera que sus vectores velocidad y aceleración tienen siempre longitud constante, probar que la curvatura de la curva es constante en todos sus puntos.
\end{ejer}

%%%%%%%%%%%%%%%%%%%%%%%%%%%%%%%%%%%%%%%%%%%
\begin{ejer}
Determinar el punto de la parábola de ecuación $y=x^2$ en el que la curvatura alcanza su valor máximo.
\end{ejer}

\begin{proof}[Solución]
Definamos la trayectoria
    \[
        \funcion{\alpha}{\R}{\R^2}{t}{(t,t^2).}
    \]
    Tenemos que
    \[
        \kappa(t)=\frac{1}{\rho(t)}\|T'(t)\|,
    \]
    para todo $t\in\R$, donde
    \begin{align*}
        T'(t)
            & = \l(\frac{\alpha'(t)}{\|\alpha'(t)\|}\r)' \\
            & = \frac{\rho(t)\alpha''(t)-\rho'(t)\alpha'(t)}{\rho^3(t)}.
    \end{align*}
    Así,
    \[
        \kappa(t)=\frac{2}{\l(\sqrt{1+4t^2}\r)^3}
    \]
    para todo $t\in\R$.
    Ahora, buscamos $t\in\R$ de modo que $\kappa$ alcance su valor máximo. Notemos que $\kappa$ es diferenciable en $\R$, por lo tanto
    \[
        \kappa'(t)=-\frac{24t}{(1+4t^2)^\frac{5}{2}},
    \]
    para todo $t\in\R$, de donde se tiene que $\kappa'(t)=0$ si y solo si $t=0$, es decir que $t=0$ es un punto crítico para $\kappa$. Finalmente, notemos que $\kappa'(t)>0$ si $t<0$ y $\kappa'(t)<0$ si $t>0$, es decir, $\kappa$ alcanza su valor máximo cuando $t=0$.
\end{proof}

%%%%%%%%%%%%%%%%%%%%%%%%%%%%%%%%%%%%%%%%%%%
\begin{ejer}
Demostrar que en una circunferencia de radio 1 la curvatura es constante.
\end{ejer}

\begin{proof}[Solución]
Consideremos la trayectoria
    \[
        \funcion{\alpha}{[0,2\pi]}{\R^2}{t}{(\cos(t),\sin(t)).}
    \]
    Se tiene que
    \[
        \rho(t)=1
        \texty
        T(t)=(-\sen(t),\cos(t)),
    \]
    para todo $t\in[0,2\pi]$, por lo tanto
    \begin{align*}
        \kappa(t)
            & = \frac{1}{\rho(t)}\|T'(t)\| \\
            & = \|(-\cos(t),-\sen(t))\| \\
            & = 1,
    \end{align*}
    para todo $t\in[0,2\pi]$, lo que implica que la curvatura es constante.
\end{proof}
    %%%%%%%%%%%%%%%%%%%%%%%%%%%%%%%%%%%%%%%%%%% JO
%%%%%%%%%%%%%%%%%%%%%%%%%%%%%%%%%%%%%%%%%%%
%%%%%%%%%%%%%%%%%%%%%%%%%%%%%%%%%%%%%%%%%%%
% Tema 5: Limites y diferenciabilidad
\item 
    Sea el campo escalar
    \[
        \funcion{f}{\R^2}{\R}{(x,y)}{xy^2.}
    \]
    Dado que $f$ es un polinomio, $f$ es continuo y se cumple que
    \[
        \Lim_{(x,y)\to (3,2)}f(x,y)=f(3,2)=12.
    \]
    Para $\epsilon=1$, determine explícitamente un $\delta>0$ tal que
    \[
    (\forall (x,y)\in\R^2)(\|(x,y)-(3,2)\|< \delta \Imp |xy^2-12|<1).
    \]

\begin{proof}    
    De la definición de norma, para cada $(x,y)\in\R^2$, se tiene que
    \[
        \|(x,y)-(3,2)\|=\sqrt{(x-3)^2+(y-2)},
    \]
    luego $\|(x,y)-(3,2)\|\geq |x-3|$ y $\|(x,y)-(3,2)\|\geq |y-2|$. Queremos acotar el término
    \[
        |xy^2-12|
    \]
    cerca del punto $(3,2)$. Con esto y dado que $(x,y)\to (3,2)$,  $|x-3|$ y $|y-2|$ deben ser términos pequeños, se tiene que para cada $(x,y)\in\R^2$, 
    \begin{equation}\label{eq:dtri}
        \begin{split}|xy^2-12|
        &=|(x-3)(y^2-4)+4x+3y^2-24|\\
        &=|(x-3)(y^2-4)+4(x-3)+3(y^2-4)|\\
        &\leq |x-3||y^2-4|+4|x-3|+3|y^2-4|\\
        &=|x-3||y-2||y+2|+4|x-3|+3|y-2||y+2|.
        \end{split}
    \end{equation}
    Así, debemos acotar cada uno de los términos anteriores. Como el $\delta$ debe ser pequeño, podemos suponer que $\delta\leq 1$. Entonces si $(x,y)\in\R^2$ tales que
    \[
        \|(x,y)-(3,2)\|<\delta\leq 1,
    \]
    se tiene que 
    \[
        |x-3|\leq \|(x,y)-(3,2)\|<\delta \yds 
        |y-2|\leq \|(x,y)-(3,2)\|<\delta.
    \]
    Ahora bien, como $|y-2|<1$, se tiene que $|y+2|<5$. Por todo esto,
    \[
        |x-3||y-2||y+2|+4|x-3|+3|y-2||y+2|
        <(\delta)( \delta) 5+4\delta +3 (\delta) 5
        =5\delta^2+19\delta. 
    \]
    De la última desigualdad y \eqref{eq:dtri}, es suficiente encontrar $\delta>0$, tal que
    \[
        5\delta^2+19\delta\leq 1
    \]
    lo qué ocurre cuando 
    \[
        0<\delta\leq \dfrac{ \sqrt{381} - 19}{10}.
    \]
    Por lo tanto, es suficiente tomar
    \[
        \delta=\dfrac{ \sqrt{381} - 19}{10}\approx 0.0519. \qedhere
    \]
\end{proof}


\end{document}
